% Introduction. Source: CLAUDE.md §1 (one line, pipeline, contributions, what we do NOT claim,
% research questions). Numbers only by reference to the chapters' generated tables (rule 11).

\chapter{Introduction}
\label{ch:intro}

Thermal design of a circuit board is a search: move chips, pour copper, place thermal vias, and
check the peak temperature after every change. An exact solve of the steady heat equation on a
four-layer board at useful resolution takes seconds, too slow to sit inside an interactive loop or
an optimiser that proposes thousands of designs. A neural surrogate answers in milliseconds but
can be wrong, and nothing in its output says when. TwinHeat combines the two so that the fast
answer can be trusted where it is used and checked where it is not, with guarantees on both sides.

\section{The pipeline}
A five-member ensemble of Fourier neural operators predicts the steady temperature rise of a
$100\times100$~mm four-layer board on a $256\times256\times4$ grid. For each proposed design an
\emph{audited adaptive conformal gate} decides whether to accept the prediction or to flag the
design to the exact solver; a random audit coin, drawn only after the design is committed and
never shown to the proposer, sends some accepted designs to the solver as well, and those audits
are the only data that move the gate's threshold. Flagged and audited designs are verified by a
warm-started conjugate-gradient solve that stops as soon as an \emph{M-matrix certificate}, a
guaranteed bound on the error in every cell, settles its question. A WebXR viewer shows the field
converging inside that guaranteed band.

\section{Contributions}
\begin{enumerate}
\item An M-matrix certificate with certified early stopping of verifications and audits: every
bound is a true upper bound on the error to the exact discrete solution, including floating-point
rounding, and a solve stops when its decision is settled rather than at a fixed tolerance
(Chapter~\ref{ch:pagp}, Appendix~\ref{app:lemmas}).
\item Validity of the gate for any proposer, a human included, with certified audit labels
(Lemma~1 and 1$'$; Chapter~\ref{ch:dl}). A per-group version was a target; it is reported as an
empirical result, not a proved one.
\item GPU certify-on-demand: one batched, warm-started launch in which each board leaves when it is
certified (Chapter~\ref{ch:pagp}).
\item Rigorous per-iterate bands in an interactive, anytime-rendered viewer, and a live system in
which the gate, the coin and certified verification run behind the user's edits
(Chapter~\ref{ch:cgvr}).
\end{enumerate}

\paragraph{What is not claimed.} Conformal bands for neural operators; online conformal inference
with intermittently revealed labels, including the transfer to all rounds that Lemma~1 restates;
audits as a source of unbiased labels; active labelling for sequential risk control; gating a
surrogate to a warm-started solver; warm starts as such; energy-form physics losses; learned
preconditioners; energy- and $A$-norm error bounds; and the inequality
$\|e\|\le\|A^{-1}\|\,\|r\|$. All of these are prior work. Learned certificates failed a
preliminary test and are not claimed. Hardware validation against fabricated boards was planned
and dropped.

\section{Research questions and structure}
The work asks whether the gate keeps accepted errors rare for any proposer and per group (RQ1), at
what cost in exact solves (RQ2), how early the certificate settles a verification (RQ3), how many
certified designs per second the GPU pipeline delivers (RQ4), and whether the calibrated signal
helps designers (RQ5; the user study was not run). Chapter~\ref{ch:dl} covers the surrogate and
the gate, Chapter~\ref{ch:pagp} the solver, the certificate and the GPU pipeline,
Chapter~\ref{ch:cgvr} the viewer and the live system, and Chapter~\ref{ch:conclusion} summarises
the evidence for each hypothesis. Every number in the report comes from a script run at a tagged
commit, and every table names the result file it was generated from.
